\documentclass[11pt]{article}
\usepackage[margin=1in]{geometry}
\usepackage{amsmath,amssymb}
\usepackage{hyperref}
\usepackage{parskip}

\title{A History of Quantum Mechanics in Experiments}
\author{}
\date{\today}

\begin{document}
\maketitle

\begin{abstract}
Quantum mechanics is, more than most physical theories, a creature of the laboratory. Its concepts---quantization, wave--particle duality, spin, entanglement, gauge topology---were forced on theorists by experiments that refused to fit classical pictures. This article traces the canonical experiments in roughly chronological order, from the ultraviolet catastrophe to satellite-based Bell tests, noting in each case what was done, by whom, and what was learned.
\end{abstract}

\section{The Old Quantum Theory (1900--1925)}

\subsection{Black-body radiation (Planck, 1900)}
Building on precision spectral measurements by Lummer, Pringsheim, Rubens, and Kurlbaum at the Physikalisch-Technische Reichsanstalt, Max Planck derived his radiation law by postulating that oscillators in the cavity walls could only exchange energy in discrete units $E = h\nu$. The constant $h$, introduced as a mathematical fix, became the cornerstone of quantum theory.

\subsection{The photoelectric effect (Hertz 1887; Einstein 1905; Millikan 1916)}
Heinrich Hertz observed in 1887 that ultraviolet light facilitates electrical discharge, and Philipp Lenard showed that the kinetic energy of ejected electrons depends on the light's frequency, not its intensity. Einstein's 1905 \emph{Annus Mirabilis} paper explained this by treating light itself as quantized into photons of energy $h\nu$. Robert Millikan, initially skeptical, spent a decade attempting to disprove Einstein's formula and instead confirmed it to high precision in 1916, simultaneously providing one of the best early measurements of $h$.

\subsection{Millikan's oil-drop experiment (1909)}
Robert Millikan and Harvey Fletcher measured the elementary charge $e$ by balancing gravity against an electric field acting on charged oil droplets. The result---that all measured charges were integer multiples of a single fundamental unit---established the quantization of electric charge.

\subsection{Rutherford scattering (Geiger--Marsden, 1909)}
Under Ernest Rutherford's direction, Hans Geiger and Ernest Marsden fired alpha particles at a thin gold foil and observed that a small fraction scattered through very large angles. Rutherford concluded in 1911 that atoms consist of a tiny dense positive nucleus surrounded by mostly empty space, overturning Thomson's plum-pudding model and setting the stage for Bohr's atom.

\subsection{Franck--Hertz experiment (1914)}
James Franck and Gustav Hertz accelerated electrons through mercury vapor and observed that the current dropped sharply at specific accelerating voltages (4.9 V for Hg). This demonstrated directly that atoms absorb energy only in discrete amounts, providing independent experimental confirmation of Bohr's quantized energy levels.

\subsection{Einstein--de Haas effect (1915)}
Albert Einstein and Wander de Haas suspended an iron cylinder inside a solenoid and observed that magnetizing it caused mechanical rotation, demonstrating that magnetization is associated with angular momentum. Their measured gyromagnetic ratio was off by a factor of two---an early hint of electron spin, though this would not be appreciated for another decade.

\subsection{Stern--Gerlach experiment (1922)}
Otto Stern and Walther Gerlach passed a beam of silver atoms through an inhomogeneous magnetic field and observed that the beam split into exactly two components rather than smearing continuously. This was originally interpreted as confirming Sommerfeld's space quantization of orbital angular momentum, but after Uhlenbeck and Goudsmit's 1925 proposal it was recognized as a direct observation of electron spin-$\tfrac{1}{2}$.

\subsection{Compton scattering (1923)}
Arthur Compton scattered X-rays off graphite and measured a wavelength shift dependent on scattering angle, exactly matching the kinematics of an elastic collision between a photon (carrying momentum $h\nu/c$) and a free electron. This was the decisive experimental evidence that light quanta carry momentum and behave as particles in collisions.

\section{Wave--Particle Duality (1927--1989)}

\subsection{Davisson--Germer experiment (1927)}
Clinton Davisson and Lester Germer at Bell Labs scattered low-energy electrons off a nickel single crystal and observed diffraction peaks consistent with de Broglie's hypothesis $\lambda = h/p$. George Paget Thomson independently demonstrated electron diffraction through thin metal foils the same year. Matter waves were real.

\subsection{The double-slit experiment with electrons (J\"onsson 1961; Tonomura 1989)}
Although Young's 1801 double-slit experiment with light is the historical archetype, the canonical quantum-mechanical version uses massive particles. Claus J\"onsson demonstrated electron interference through multiple slits in 1961. Akira Tonomura and collaborators at Hitachi performed the definitive single-electron version in 1989, recording the gradual buildup of an interference pattern from individual electron detection events---each electron interfering with itself.

\subsection{Colella--Overhauser--Werner neutron interferometry (1975)}
Roberto Colella, Albert Overhauser, and Samuel Werner used a silicon-crystal neutron interferometer to measure the gravitationally induced phase shift on a neutron wavefunction whose two arms were at different heights in Earth's gravitational field. This was the first experiment to combine quantum mechanics and gravity in a single coherent measurement.

\section{Spin Resonance and Precision QED}

\subsection{Rabi molecular beam resonance (1938)}
Isidor Rabi developed the molecular beam magnetic resonance method, driving transitions between nuclear spin states with radiofrequency fields and detecting them via beam deflection. This technique became the foundation of NMR, atomic clocks, and modern precision spectroscopy.

\subsection{The Lamb shift (1947)}
Willis Lamb and Robert Retherford used microwave techniques developed during the war to measure a small splitting between the $2S_{1/2}$ and $2P_{1/2}$ levels of hydrogen, which Dirac's equation predicted to be degenerate. The shift, about 1058 MHz, could only be explained by the interaction of the electron with vacuum fluctuations of the electromagnetic field, and it directly motivated the development of renormalized QED by Bethe, Schwinger, Feynman, Tomonaga, and Dyson.

\subsection{Anomalous magnetic moment of the electron ($g-2$)}
Beginning with Kusch and Foley's 1947 measurement showing the electron's $g$-factor differs from Dirac's prediction of exactly 2, successive Penning-trap experiments---culminating in the work of Hans Dehmelt and later Gerald Gabrielse---have measured $g-2$ to better than one part in $10^{12}$. Agreement with QED calculations remains the most stringent test of any physical theory.

\section{Gauge Topology: The Aharonov--Bohm Effect}

In 1959 Yakir Aharonov and David Bohm pointed out that a charged particle traveling around a region containing magnetic flux---but through which the particle itself never passes---should acquire a measurable phase shift proportional to the enclosed flux. The effect implies that the electromagnetic potentials $A_\mu$, not just the field strengths $F_{\mu\nu}$, carry physical content, and that the topology of the configuration space is fundamental.

Robert Chambers reported the first experimental confirmation in 1960 using electron biprism interferometry around a magnetized whisker. Decisive confirmation came from Akira Tonomura and collaborators in 1986, who used superconducting toroidal samples to completely shield the magnetic flux from the electron paths, ruling out residual field leakage as a possible classical explanation. The Aharonov--Bohm effect remains the cleanest experimental statement that gauge structure is physical and that quantum mechanics responds to the topology, not merely the local geometry, of the spaces particles move through.

\section{Foundations: Entanglement and Nonlocality}

\subsection{The EPR argument (1935) and Bell's theorem (1964)}
Einstein, Podolsky, and Rosen argued that quantum mechanics must be incomplete because it predicts perfect correlations between spatially separated measurements that, under local-realist assumptions, would require pre-existing ``elements of reality.'' John Bell showed in 1964 that any local hidden-variable theory must satisfy an inequality that quantum mechanics violates, converting a philosophical dispute into an experimental question.

\subsection{Early Bell tests (Freedman--Clauser 1972; Aspect 1982)}
Stuart Freedman and John Clauser performed the first experimental Bell test in 1972 using polarization-entangled photons from a calcium atomic cascade, finding a violation. Alain Aspect, Philippe Grangier, and G\'erard Roger refined the experiment in 1981--82, including a version with time-varying polarization analyzers that closed the locality loophole by switching settings while the photons were in flight.

\subsection{Hong--Ou--Mandel interference (1987)}
Chung Ki Hong, Zhe-Yu Ou, and Leonard Mandel demonstrated that two indistinguishable photons arriving simultaneously at the input ports of a beam splitter always exit together, never separately. This two-photon bunching effect has no classical analogue and is now a workhorse for measuring photon indistinguishability in quantum optics and quantum information.

\subsection{Quantum eraser and delayed choice (Kim et al.\ 1999)}
John Wheeler proposed a delayed-choice gedankenexperiment in 1978 in which the decision to measure ``which path'' is made after a particle has traversed an interferometer. Yoon-Ho Kim and collaborators implemented a delayed-choice quantum eraser in 1999 using entangled photon pairs, showing that interference patterns appear or disappear depending on whether which-path information is, in principle, available---even when that determination is made after detection of the signal photon.

\subsection{Loophole-free Bell tests (2015)}
Three groups---Ronald Hanson's at Delft, Lynden Shalm's at NIST, and Anton Zeilinger's in Vienna---reported Bell inequality violations in 2015 that simultaneously closed the locality, detection, and freedom-of-choice loopholes. Hanson's experiment used entangled NV-center electron spins separated by 1.3 km on the Delft campus. These results constitute the definitive empirical refutation of local realism.

\subsection{The Micius satellite experiments (2017)}
Jian-Wei Pan's group used the Chinese Micius satellite to distribute polarization-entangled photon pairs to ground stations separated by up to 1203 km, observing Bell inequality violations at distances unreachable by terrestrial fiber. Subsequent experiments demonstrated satellite-to-ground quantum key distribution and an intercontinental quantum-secured video conference between Beijing and Vienna, establishing the feasibility of a global quantum communication network.

\subsection{Leggett--Garg tests}
Anthony Leggett and Anupam Garg proposed in 1985 a temporal analogue of Bell's inequality testing ``macroscopic realism''---the assumption that a macroscopic system possesses definite properties independent of measurement. Experimental violations using superconducting qubits, photons, and nuclear spins have placed increasingly stringent bounds on macrorealistic alternatives to quantum mechanics.

\section{Macroscopic Quantum Phenomena}

\subsection{Superconductivity and the Josephson effect}
Heike Kamerlingh Onnes discovered superconductivity in mercury in 1911. The microscopic BCS theory came in 1957, and in 1962 Brian Josephson predicted that Cooper pairs should tunnel coherently between two superconductors separated by a thin insulating barrier, producing a DC current at zero voltage and an AC current at finite voltage. Philip Anderson and John Rowell confirmed the effect experimentally in 1963. Josephson junctions and SQUIDs are now the basis of superconducting qubits and the most sensitive magnetometers ever built.

\subsection{The quantum Hall effects (1980, 1982)}
Klaus von Klitzing discovered in 1980 that the Hall conductance of a two-dimensional electron gas in a strong magnetic field is quantized in exact integer multiples of $e^2/h$, independent of sample geometry or disorder. Daniel Tsui, Horst St\"ormer, and Arthur Gossard discovered the fractional quantum Hall effect in 1982, with Robert Laughlin providing the theoretical explanation in terms of a strongly correlated incompressible quantum fluid. Both effects are topological in origin and gave the first experimental access to states of matter classified by global rather than local order parameters.

\subsection{Bose--Einstein condensation (1995)}
Eric Cornell and Carl Wieman at JILA produced the first dilute-gas BEC in rubidium-87 in June 1995, with Wolfgang Ketterle at MIT achieving condensation in sodium a few months later. Decades after Einstein's 1925 prediction, matter-wave condensation became a laboratory reality, opening the field of ultracold atomic physics.

\section{Closing Remarks}

The arc from Planck's cavity radiation formula to satellite-borne Bell tests is striking for how reluctantly each new conceptual commitment was made. Quantization, wave--particle duality, spin, entanglement, and gauge topology were not arrived at by deduction from first principles but were extracted, often grudgingly, from experiments that classical physics could not accommodate. What unites the canon is that in each case the experimental result preceded---and forced---the theoretical reorganization. The fact that quantum mechanics now admits engineering applications, from MRI to superconducting qubits to satellite QKD, is in some sense incidental to its primary achievement: a century of laboratory work has established that the world is, at its foundations, irreducibly quantum.

\end{document}
